%% ========================================================================
%% Chapter 9: Bash Programming
%% ========================================================================

\chapter{Bash Programming}
\label{ch:bash-programming}

This chapter covers the fundamentals of Bash programming: creating scripts,
managing variables, writing control structures, defining functions,
processing arguments, and debugging. The concepts are compared with Java so
that students who already know programming can build understanding more
quickly.

\begin{examplebox}[Bash as a programming language]
Imagine that you have already written Java programs. Java has variables,
conditions, loops, and functions. Bash does too, but with one important
difference: Bash executes \textit{operating system commands} as part of the
program logic. A Bash script is \textbf{a list of terminal commands stored in
a file}, enriched with control logic so it can run automatically without
human involvement.
\end{examplebox}

%%---------------------------------------------------------------------
\section{Bash Scripting Basics}
\label{sec:creating-scripts}

\subsection*{How to Create and Edit Script Files}

This chapter uses two ways to create files: \textbf{nano} for interactive
editing, and a \textbf{heredoc} for creating files directly from the
terminal.

\textbf{Using nano:}
\begin{enumerate}
  \item Open or create a file: \texttt{nano file-name.sh}
  \item Type the contents directly; the cursor is immediately active and
    there is no special mode.
  \item Important keys:
  \begin{itemize}
    \item \keystroke{Ctrl+O} then \keystroke{Enter} to save the file
    \item \keystroke{Ctrl+X} to exit nano
    \item \keystroke{Ctrl+W} to search for text in the file
  \end{itemize}
\end{enumerate}

\begin{notebox}
At the bottom of the nano screen, you will always see a list of active
shortcuts. The symbol \texttt{\^{}} means \texttt{Ctrl}, so
\texttt{\^{}O} = \keystroke{Ctrl+O}.
\end{notebox}

\textbf{Using a heredoc} (for short scripts or automation):
\begin{lstlisting}[language=bash, caption={Creating a file with a heredoc}]
cat <<'EOF' > file-name.sh
#!/bin/bash
echo "Hello from the script"
EOF
\end{lstlisting}

\subsection*{Basic Structure and the Shebang Line}

Every Bash script starts with a \textit{shebang line}:

\begin{lstlisting}[language=bash, caption={Basic structure of a Bash script}]
#!/bin/bash
# Comment: explain what this script does

echo "Script is running"
\end{lstlisting}

The \texttt{\#!} line tells the Linux kernel which interpreter to use. Two
common forms are:

\begin{itemize}
  \item \texttt{\#!/bin/bash}: the absolute path to Bash. This is the most
    common form on Linux, simple and direct.
  \item \texttt{\#!/usr/bin/env bash}: asks the \texttt{env} command to find
    Bash in \var{PATH}. This is more portable across operating systems.
\end{itemize}

\begin{center}
\begin{tabular}{lll}
\hline
\textbf{Aspect} & \textbf{\#!/bin/bash} & \textbf{\#!/usr/bin/env bash} \\
\hline
Bash location & Fixed at \texttt{/bin/bash} & Searched in \texttt{PATH} \\
Portability & Standard Linux & Linux and macOS \\
Recommendation & This lab & Cross-platform scripts \\
\hline
\end{tabular}
\end{center}

\begin{notebox}
The shebang only works if the file has execute permission (\texttt{chmod +x})
and is run directly (\texttt{./script.sh}). If you call
\texttt{bash script.sh}, the shebang is ignored because the interpreter has
already been specified explicitly.
\end{notebox}

After creating the script, give it execute permission and run it:

\begin{lstlisting}[language=bash, caption={Granting permission and running a script}]
chmod +x file-name.sh
./file-name.sh
\end{lstlisting}

\subsection*{Lab 7.1 First Script: System Report}

Objective: create, save, and run your first Bash script.

\begin{enumerate}
  \item Create the lab workspace:
\begin{lstlisting}[language=bash, caption={Preparing the workspace}]
mkdir -p ~/lab-os/chapter07/{scripts,logs,data}
cd ~/lab-os/chapter07/scripts
\end{lstlisting}

  \item Create a script with nano:
\begin{lstlisting}[language=bash, caption={Opening nano}]
nano system-report.sh
\end{lstlisting}

  \item Type the following content, save it (\keystroke{Ctrl+O}
    \keystroke{Enter}), then exit (\keystroke{Ctrl+X}):
\begin{lstlisting}[language=bash, caption={Contents of system-report.sh}]
#!/bin/bash
# Script: system-report.sh

echo "================================"
echo "  SYSTEM REPORT"
echo "================================"
echo "Date      : $(date '+%A, %d %B %Y')"
echo "Time      : $(date '+%H:%M:%S')"
echo "Hostname  : $(hostname)"
echo "User      : $(whoami)"
echo "CPU cores : $(nproc)"
echo "Free RAM  : $(free -h | awk '/^Mem/ {print $4}')"
echo "Disk /    : $(df -h / | awk 'NR==2 {print $5}') used"
echo "================================"
\end{lstlisting}

  \item Grant permission and run it:
\begin{lstlisting}[language=bash, caption={Running the script}]
chmod +x system-report.sh
./system-report.sh
\end{lstlisting}
\end{enumerate}

\begin{exercisebox}[Exercise 7.1]
Modify \texttt{system-report.sh} so that it saves the output to a file named
\texttt{report-YYYY-MM-DD.txt} while also displaying it in the terminal.
Hint: use \texttt{tee}, which was introduced in the previous chapter.
\end{exercisebox}

%%---------------------------------------------------------------------
\section{Variables and Positional Parameters}
\label{sec:variables-parameters}

\subsection*{Variables in Bash vs Java}

In Java: \texttt{String name = "Budi";}. In Bash, the rules are different:

\begin{itemize}
  \item \textbf{Declaration}: \texttt{NAME="Budi"} with no spaces around
    \texttt{=}
  \item \textbf{Accessing the value}: \texttt{\$NAME} or
    \texttt{\$\{NAME\}}
  \item \textbf{Data types}: none. All Bash variables are strings by
    default; arithmetic requires special syntax.
\end{itemize}

\begin{warningbox}
\texttt{NAME = "Budi"} (with spaces) causes an error because Bash thinks
\texttt{NAME} is a command name. Always write assignments without spaces.
\end{warningbox}

The special variables most often used in scripting are:

\begin{center}
\begin{tabular}{ll}
\hline
\textbf{Variable} & \textbf{Value} \\
\hline
\texttt{\$0} & Name of the currently running script \\
\texttt{\$1}, \texttt{\$2}, \ldots & First argument, second argument, etc. \\
\texttt{\$\#} & Number of arguments provided \\
\texttt{\$@} & All arguments as a separated list \\
\texttt{\$?} & Exit code of the last command (0 = success) \\
\texttt{\$\$} & PID of the current script \\
\hline
\end{tabular}
\end{center}

\subsection*{Positional Parameters and Default Values}

Just like a Java method that accepts arguments
(\texttt{void function(String a, String b)}), a Bash script accepts arguments
from the command line. Arguments are read with \texttt{\$1}, \texttt{\$2},
and so on.

The form \texttt{\$\{1:-"default"\}} reads \texttt{\$1}; if it is empty,
use a default value instead. This is equivalent to
\texttt{args.length > 0 ? args[0] : "default"} in Java.

\subsection*{Arithmetic and Command Substitution}

Bash does not calculate arithmetic expressions directly. Use
\texttt{\$(( ))}:

\begin{lstlisting}[language=bash, caption={Arithmetic and command substitution}]
A=10
B=3

echo "Sum       : $((A + B))"
echo "Modulo    : $((A % B))"
echo "Division  : $((A / B))"   # Integer result, just like int in Java

# Command substitution: store output in a variable
DATE=$(date '+%F')
FILE_COUNT=$(ls | wc -l)
echo "Date: $DATE, Files: $FILE_COUNT"
\end{lstlisting}

\begin{notebox}
In Java, \texttt{10 / 3} with type \texttt{int} produces \texttt{3}. In
Bash, \texttt{\$((10 / 3))} also produces \texttt{3}. Use \texttt{bc} if
you need decimal precision.
\end{notebox}

\subsection*{Lab 7.2 System Info Script with Arguments}

Objective: practice variables, command substitution, positional parameters,
and default values.

\begin{enumerate}
  \item Create the script:
\begin{lstlisting}[language=bash, caption={Creating system-info.sh}]
nano ~/lab-os/chapter07/scripts/system-info.sh
\end{lstlisting}

  \item Type the following contents:
\begin{lstlisting}[language=bash, caption={Contents of system-info.sh}]
#!/bin/bash
# Usage: ./system-info.sh [admin-name] [disk-threshold-percent]

ADMIN=${1:-"Unknown"}
THRESHOLD=${2:-80}
DATE=$(date '+%F %T')
DISK_PERCENT=$(df / | awk 'NR==2 {print $5}' | tr -d '%')

echo "Admin     : $ADMIN"
echo "Date      : $DATE"
echo "Disk /    : ${DISK_PERCENT}% used"
echo "Threshold : ${THRESHOLD}%"

if [ "$DISK_PERCENT" -gt "$THRESHOLD" ]; then
    echo "STATUS    : WARNING - disk usage exceeds threshold!"
else
    REMAINING=$((THRESHOLD - DISK_PERCENT))
    echo "STATUS    : Normal (${REMAINING}% tolerance remaining)"
fi
\end{lstlisting}

  \item Save it, grant permission, and test it with different argument
    combinations:
\begin{lstlisting}[language=bash, caption={Testing the script}]
chmod +x ~/lab-os/chapter07/scripts/system-info.sh
./system-info.sh
./system-info.sh "Dian" 50
./system-info.sh "Dian" 10
\end{lstlisting}
\end{enumerate}

\begin{exercisebox}[Exercise 7.2]
Create a script named \texttt{calculator.sh} that accepts three arguments:
\texttt{<number1> <operator> <number2>} where the operator is one of
\texttt{+}, \texttt{-}, \texttt{*}, or \texttt{/}. Example:
\texttt{./calculator.sh 20 + 5} produces \texttt{25}. Use \texttt{case} to
choose the operation, and validate incomplete arguments.
\end{exercisebox}

%%---------------------------------------------------------------------
\section{Control Structures}
\label{sec:control-structures}

Control structures in Bash follow the same concepts as Java, but the syntax is
different.

\begin{center}
\begin{tabular}{lll}
\hline
\textbf{Concept} & \textbf{Java} & \textbf{Bash} \\
\hline
If condition & \texttt{if (x > 5) \{} & \texttt{if [ \$x -gt 5 ]; then} \\
Code block & curly braces & \texttt{then ... fi} \\
Else if & \texttt{else if (...) \{} & \texttt{elif ...; then} \\
Not equal & \texttt{!=} & \texttt{!=} (string), \texttt{-ne} (number) \\
And / Or & \texttt{\&\&} / \texttt{||} & \texttt{\&\&} / \texttt{||} \\
Switch & \texttt{switch (x) \{ case "a": \}} & \texttt{case \$x in "a") ... ;;} \\
\hline
\end{tabular}
\end{center}

\subsection*{Comparison Operators}

\textbf{Numbers:} \texttt{-eq}, \texttt{-ne}, \texttt{-gt}, \texttt{-lt},
\texttt{-ge}, \texttt{-le}

\textbf{Strings:} \texttt{=} or \texttt{==} (equal), \texttt{!=} (not
equal), \texttt{-z} (empty), \texttt{-n} (not empty)

\textbf{Files:} \texttt{-f} (file exists), \texttt{-d} (directory exists),
\texttt{-r} (readable), \texttt{-x} (executable)

\subsection*{\texttt{for} and \texttt{while} Loops}

\begin{lstlisting}[language=bash, caption={for and while loops}]
# for loop over a list (similar to Java for-each)
for MONTH in January February March; do
    echo "Month: $MONTH"
done

# C-style for loop (similar to Java)
for ((I=1; I<=5; I++)); do
    echo "Number: $I"
done

# loop over files
for FILE in ~/lab-os/chapter07/scripts/*.sh; do
    echo "Script: $(basename "$FILE")"
done

# while loop while condition is true
COUNTER=1
while [ $COUNTER -le 3 ]; do
    echo "Iteration $COUNTER"
    COUNTER=$((COUNTER + 1))
done
\end{lstlisting}

\subsection*{The \texttt{case} Structure}

\texttt{case} is equivalent to \texttt{switch} in Java. Each block ends with
\texttt{;;} (not \texttt{break}). The pattern \texttt{*)} is equivalent to
\texttt{default:}.

\begin{notebox}
In Java, forgetting to write \texttt{break} causes \textit{fall-through}. In
Bash, \texttt{;;} already stops execution of the current block, so there is no
fall-through by default.
\end{notebox}

\subsection*{Lab 7.3 Grading Script and Interactive Menu}

Objective: combine \texttt{if}, \texttt{for}, \texttt{while}, and
\texttt{case} in one session.

\begin{enumerate}
  \item Create a grading script (using \texttt{if} and \texttt{for}):
\begin{lstlisting}[language=bash, caption={Creating batch-grading.sh}]
nano ~/lab-os/chapter07/scripts/batch-grading.sh
\end{lstlisting}

  \item Type the following contents:
\begin{lstlisting}[language=bash, caption={Contents of batch-grading.sh}]
#!/bin/bash
# Script: batch-grading.sh
# Process a list of student scores

STUDENTS=("Andi:92" "Budi:73" "Citra:55" "Deni:80" "Eka:45")

echo "=== GRADING RESULTS ==="
for ENTRY in "${STUDENTS[@]}"; do
    NAME=$(echo "$ENTRY" | cut -d: -f1)
    SCORE=$(echo "$ENTRY" | cut -d: -f2)

    if [ "$SCORE" -ge 85 ]; then
        GRADE="A"
    elif [ "$SCORE" -ge 75 ]; then
        GRADE="B"
    elif [ "$SCORE" -ge 65 ]; then
        GRADE="C"
    elif [ "$SCORE" -ge 55 ]; then
        GRADE="D"
    else
        GRADE="E"
    fi

    printf "%-10s %3d  %s\n" "$NAME" "$SCORE" "$GRADE"
done
echo "======================="
\end{lstlisting}

  \item Save it, grant permission, and run it:
\begin{lstlisting}[language=bash, caption={Running the grading script}]
chmod +x ~/lab-os/chapter07/scripts/batch-grading.sh
./batch-grading.sh
\end{lstlisting}

  \item Create an interactive menu script (\texttt{while} + \texttt{case}):
\begin{lstlisting}[language=bash, caption={Creating system-menu.sh}]
nano ~/lab-os/chapter07/scripts/system-menu.sh
\end{lstlisting}

  \item Type the following contents:
\begin{lstlisting}[language=bash, caption={Contents of system-menu.sh}]
#!/bin/bash
# Interactive system monitoring menu

while true; do
    echo ""
    echo "===== MONITOR MENU ====="
    echo "1) Disk info"
    echo "2) Memory info"
    echo "3) Top processes"
    echo "4) Exit"
    echo -n "Choose [1-4]: "
    read CHOICE

    case $CHOICE in
        1) df -h ;;
        2) free -h ;;
        3) ps aux --sort=-%cpu | head -6 ;;
        4) echo "Goodbye!"; exit 0 ;;
        *) echo "Invalid choice." ;;
    esac
done
\end{lstlisting}

  \item Grant permission and run it, trying each option:
\begin{lstlisting}[language=bash, caption={Running the menu}]
chmod +x ~/lab-os/chapter07/scripts/system-menu.sh
./system-menu.sh
\end{lstlisting}
\end{enumerate}

\begin{exercisebox}[Exercise 7.3]
Add a summary section to the bottom of \texttt{batch-grading.sh} that shows
the number of students in each grade (A, B, C, D, E) using a second
\texttt{for} loop that iterates over the \texttt{STUDENTS} array.
\end{exercisebox}

%%---------------------------------------------------------------------
\section{Functions in Shell Scripts}
\label{sec:functions-shell-scripts}

\subsection*{Functions in Bash vs Java}

In Java: \texttt{void methodName(String p1, String p2) \{ ... \}}. In Bash,
functions are defined without an explicit parameter list; arguments are passed
like command-line arguments and read with \texttt{\$1}, \texttt{\$2}, and so
on.

\begin{center}
\begin{tabular}{lll}
\hline
\textbf{Concept} & \textbf{Java} & \textbf{Bash} \\
\hline
Definition & \texttt{void f(String a, int b)} & \texttt{f() \{ ... \}} \\
Call & \texttt{f("x", 2)} & \texttt{f x 2} \\
Arguments & \texttt{a}, \texttt{b} & \texttt{\$1}, \texttt{\$2} \\
Return value & \texttt{return value} & \texttt{echo value} (via stdout) \\
Return status & exception & \texttt{return 0--255} \\
Local variables & local by default & must write \texttt{local} \\
\hline
\end{tabular}
\end{center}

\begin{notebox}
In Bash, \texttt{return} only returns a status code (0--255). To return a
value such as a string, use \texttt{echo} inside the function, then capture
it with substitution: \texttt{RESULT=\$(function\_name arg)}. In Java,
variables inside methods are automatically local. In Bash, without
\texttt{local}, variables leak into the global scope.
\end{notebox}

\begin{lstlisting}[language=bash, caption={Defining and calling functions}]
#!/bin/bash

# Function returns a value via stdout
calculate_area() {
    local LENGTH=$1
    local WIDTH=$2
    echo $((LENGTH * WIDTH))
}

# Capture the value with command substitution
AREA=$(calculate_area 8 5)
echo "Area: $AREA cm2"

# Function returns success/failure status
divide() {
    local A=$1 B=$2
    if [ "$B" -eq 0 ]; then
        echo "Error: division by zero" >&2
        return 1
    fi
    echo $((A / B))
}

RESULT=$(divide 20 4) && echo "Result: $RESULT"
divide 10 0 || echo "Exit code: $?"
\end{lstlisting}

\subsection*{Lab 7.4 Validation Function Library}

Objective: create a function library file that can be loaded
(\texttt{source}) by other scripts, similar to \textit{import} in Java.

\begin{enumerate}
  \item Create the library file:
\begin{lstlisting}[language=bash, caption={Creating lib-validation.sh}]
nano ~/lab-os/chapter07/scripts/lib-validation.sh
\end{lstlisting}

  \item Type the following contents:
\begin{lstlisting}[language=bash, caption={Contents of lib-validation.sh}]
#!/bin/bash
# lib-validation.sh - Validation function library

is_number() {
    [[ "$1" =~ ^[0-9]+$ ]]
}

is_readable_file() {
    [ -f "$1" ] && [ -r "$1" ]
}

error_exit() {
    echo "ERROR: $1" >&2
    exit 1
}

info()    { echo "[INFO] $1"; }
success() { echo "[OK]   $1"; }
\end{lstlisting}

  \item Create a script that uses the library:
\begin{lstlisting}[language=bash, caption={Creating use-library.sh}]
nano ~/lab-os/chapter07/scripts/use-library.sh
\end{lstlisting}

  \item Type the following contents:
\begin{lstlisting}[language=bash, caption={Contents of use-library.sh}]
#!/bin/bash
# Load the library (similar to import in Java)
source ~/lab-os/chapter07/scripts/lib-validation.sh

NUMBER=$1
FILE=$2

[ -z "$NUMBER" ] || [ -z "$FILE" ] && \
    error_exit "Usage: $0 <number> <file-path>"

if is_number "$NUMBER"; then
    success "Input '$NUMBER' is a valid number"
else
    error_exit "'$NUMBER' is not a number"
fi

if is_readable_file "$FILE"; then
    success "File '$FILE' is readable"
    info "Line count: $(wc -l < "$FILE")"
else
    error_exit "File '$FILE' does not exist or is not readable"
fi
\end{lstlisting}

  \item Grant permission and test every scenario:
\begin{lstlisting}[language=bash, caption={Testing all scenarios}]
chmod +x ~/lab-os/chapter07/scripts/use-library.sh
./use-library.sh 42 /etc/hostname
./use-library.sh abc /etc/hostname
./use-library.sh 42 /does-not-exist.txt
./use-library.sh
\end{lstlisting}
\end{enumerate}

\begin{exercisebox}[Exercise 7.4]
Add a \texttt{confirm()} function to \texttt{lib-validation.sh}. This
function should display a question, read Y/N input from the user, return
\texttt{0} for Y and \texttt{1} for N. Create a demo script that calls this
function before deleting a file.
\end{exercisebox}

%%---------------------------------------------------------------------
\section{Command-Line Argument Processing}
\label{sec:argument-processing}

Good scripts process arguments flexibly: validate the number of arguments, and
handle options such as \texttt{-v} or \texttt{-o file}.

\texttt{getopts} is Bash's standard way to process options, similar to how
Unix commands such as \texttt{ls -la} or \texttt{grep -i} work:

\begin{lstlisting}[language=bash, caption={Basic getopts pattern}]
#!/bin/bash
VERBOSE=false
OUTPUT_FILE=""

while getopts "vo:" OPTION; do
    case $OPTION in
        v) VERBOSE=true ;;
        o) OUTPUT_FILE="$OPTARG" ;;
        *) echo "Usage: $0 [-v] [-o file]"; exit 1 ;;
    esac
done
shift $((OPTIND - 1))   # Shift so that $1 = first non-option argument

echo "Verbose : $VERBOSE"
echo "Output  : ${OUTPUT_FILE:-none}"
echo "Remain  : $@"
\end{lstlisting}

In \texttt{"vo:"}, the \texttt{:} after a letter means that option requires
a value (\texttt{-o file-name}), not just a flag (\texttt{-v}).

\subsection*{Lab 7.5 Backup Script with Options}

Objective: create a script that processes \texttt{getopts} options and
integrates all the earlier concepts.

\begin{enumerate}
  \item Create the script:
\begin{lstlisting}[language=bash, caption={Creating backup-data.sh}]
nano ~/lab-os/chapter07/scripts/backup-data.sh
\end{lstlisting}

  \item Type the following contents:
\begin{lstlisting}[language=bash, caption={Contents of backup-data.sh}]
#!/bin/bash
# Usage: ./backup-data.sh [-v] [-c] [-l logfile] <source> <destination>

VERBOSE=false
COMPRESS=false
LOG_FILE=""

while getopts "vcl:" OPTION; do
    case $OPTION in
        v) VERBOSE=true ;;
        c) COMPRESS=true ;;
        l) LOG_FILE="$OPTARG" ;;
        *) echo "Usage: $0 [-v] [-c] [-l logfile] <source> <destination>"
           exit 1 ;;
    esac
done
shift $((OPTIND - 1))

SOURCE=$1
DESTINATION=$2

log() {
    local MSG="[$(date '+%T')] $1"
    echo "$MSG"
    [ -n "$LOG_FILE" ] && echo "$MSG" >> "$LOG_FILE"
}

[ -z "$SOURCE" ] || [ -z "$DESTINATION" ] && {
    echo "ERROR: source and destination are required"; exit 1; }

[ ! -d "$SOURCE" ] && { log "ERROR: $SOURCE does not exist"; exit 2; }

mkdir -p "$DESTINATION"
DATE=$(date '+%F-%H%M%S')
[ "$VERBOSE" = true ] && log "Source: $SOURCE | Destination: $DESTINATION"

if [ "$COMPRESS" = true ]; then
    ARCHIVE="$DESTINATION/backup-$(basename "$SOURCE")-$DATE.tar.gz"
    tar -czf "$ARCHIVE" -C "$(dirname "$SOURCE")" "$(basename "$SOURCE")"
    log "Archive: $ARCHIVE ($(du -sh "$ARCHIVE" | cut -f1))"
else
    cp -r "$SOURCE" "$DESTINATION/backup-$(basename "$SOURCE")-$DATE"
    log "Backup completed."
fi
\end{lstlisting}

  \item Grant permission and test it:
\begin{lstlisting}[language=bash, caption={Testing the backup script}]
chmod +x ~/lab-os/chapter07/scripts/backup-data.sh
cd ~/lab-os/chapter07/scripts

# Without options
./backup-data.sh ~/lab-os/chapter07/data ~/lab-os/chapter07/logs

# With verbose and compression + logging to file
./backup-data.sh -v -c -l ../logs/backup.log \
  ~/lab-os/chapter07/data ~/lab-os/chapter07/logs

cat ../logs/backup.log
\end{lstlisting}
\end{enumerate}

%%---------------------------------------------------------------------
\section{Debugging Shell Scripts}
\label{sec:debugging-script}

Debugging in Bash is different from Java, which has interactive debuggers.
Bash provides tracing mechanisms to show each line before it is executed.

\begin{center}
\begin{tabular}{lll}
\hline
\textbf{Flag} & \textbf{Meaning} & \textbf{How to enable it} \\
\hline
\texttt{-x} & Show each command before it runs & \texttt{bash -x script.sh} \\
\texttt{-e} & Exit when a command fails & \texttt{set -e} \\
\texttt{-u} & Exit on undefined variables & \texttt{set -u} \\
\texttt{-n} & Check syntax without executing & \texttt{bash -n script.sh} \\
\hline
\end{tabular}
\end{center}

The combination \texttt{set -euo pipefail} at the top of a script catches many
common errors automatically. Add \texttt{trap} for cleanup:

\begin{lstlisting}[language=bash, caption={Defensive script pattern}]
#!/bin/bash
set -euo pipefail

cleanup() { echo "[CLEANUP] done."; }
trap cleanup EXIT      # cleanup always runs when the script ends

DEBUG=${DEBUG:-false}
debug() { [ "$DEBUG" = true ] && echo "[DEBUG] $*" >&2; }

debug "Script started with PID $$"
echo "Script is running normally"
\end{lstlisting}

\begin{notebox}
\texttt{trap cleanup EXIT} is similar to \texttt{finally} in Java: it ensures
cleanup runs even if the script fails partway through.
\end{notebox}

\subsection*{Lab 7.6 Debugging a Script}

Objective: use debugging techniques to analyze and fix a script.

\begin{enumerate}
  \item Create the script to analyze:
\begin{lstlisting}[language=bash, caption={Creating debug-practice.sh}]
nano ~/lab-os/chapter07/scripts/debug-practice.sh
\end{lstlisting}

  \item Type the following contents:
\begin{lstlisting}[language=bash, caption={Contents of debug-practice.sh}]
#!/bin/bash
# Script: debug-practice.sh
# Usage: ./debug-practice.sh <directory> <limit-MB>

DIRECTORY=$1
LIMIT=$2

if [ $# -ne 2 ]; then
    echo "Usage: $0 <directory> <limit-MB>"
    exit 1
fi

SIZE=$(du -sm "$DIRECTORY" | cut -f1)

echo "Directory : $DIRECTORY"
echo "Size      : ${SIZE} MB"
echo "Limit     : ${LIMIT} MB"

if [ "$SIZE" -gt "$LIMIT" ]; then
    echo "WARNING: Size exceeds the limit!"
    echo "Excess: $((SIZE - LIMIT)) MB"
else
    echo "Status: Normal (remaining: $((LIMIT - SIZE)) MB)"
fi
\end{lstlisting}

  \item Check the syntax, then run it with tracing:
\begin{lstlisting}[language=bash, caption={Analyzing the script with debugging}]
chmod +x ~/lab-os/chapter07/scripts/debug-practice.sh
bash -n debug-practice.sh && echo "Syntax OK"
bash -x debug-practice.sh /etc 10
./debug-practice.sh /var 50
./debug-practice.sh
\end{lstlisting}
\end{enumerate}

\begin{exercisebox}[Exercise 7.5]
The script \texttt{debug-practice.sh} does not handle missing directories.
Fix it by adding:
\begin{itemize}
  \item \texttt{set -e} on the second line
  \item A \texttt{-d "\$DIRECTORY"} check before calling \texttt{du}
  \item An informative error message if the directory is not found
\end{itemize}
Test it with a directory that does not exist.
\end{exercisebox}

%%---------------------------------------------------------------------
\section{Best Practices}
\label{sec:best-practices}

Good Bash scripts follow these principles:

\begin{enumerate}
  \item \textbf{Always write a shebang} (\texttt{\#!/bin/bash}).
  \item \textbf{Use \texttt{set -euo pipefail}} at the start of the script.
  \item \textbf{Quote all variables}: \texttt{"\$NAME"}, not
    \texttt{\$NAME}, to prevent problems when values contain spaces.
  \item \textbf{Use \texttt{local}} inside functions to prevent variables
    from leaking into the global scope.
  \item \textbf{Validate all input} before processing it.
  \item \textbf{Use descriptive variable names}: \texttt{BACKUP\_DIRECTORY}
    is clearer than \texttt{D}.
  \item \textbf{Use \texttt{trap}} for cleanup when the script exits.
  \item \textbf{Separate configuration from logic}: place configuration
    variables at the top of the script.
\end{enumerate}

Professional script template that applies all of the principles above:

\begin{longlisting}[language=bash, caption={Professional Bash script template}]
#!/bin/bash
# ====================================================
# Script: script-name.sh
# Description: Brief explanation of this script
# Usage: ./script-name.sh [-v] [-h] <arguments>
# ====================================================

set -euo pipefail

SCRIPT_NAME=$(basename "$0")
VERBOSE=false

usage() {
    echo "Usage: $SCRIPT_NAME [-v] [-h] <arguments>"
    echo "  -v  Verbose mode"
    echo "  -h  Show help"
    exit 0
}

log()         { echo "[$SCRIPT_NAME] $*"; }
log_verbose() { [ "$VERBOSE" = true ] && log "$*"; }
error_exit()  { echo "ERROR: $*" >&2; exit 1; }
cleanup()     { log_verbose "Cleanup finished."; }

trap cleanup EXIT

while getopts "vh" OPTION; do
    case $OPTION in
        v) VERBOSE=true ;;
        h) usage ;;
        *) error_exit "Unknown option. Use -h." ;;
    esac
done
shift $((OPTIND - 1))

[ $# -lt 1 ] && error_exit "Incomplete arguments. Use -h."

log "Script started"
log_verbose "Arguments: $*"

# Main logic goes here

log "Done."
\end{longlisting}

%%---------------------------------------------------------------------
\section{Lab Assignments}
\label{sec:lab-assignment-bash-prog}

\begin{warningbox}
Do the assignments in \texttt{\$HOME/lab-os/chapter07/}. Put scripts in
\texttt{scripts/} and output files in \texttt{logs/}.
\end{warningbox}

\subsection*{Assignment 1 Class Attendance Script}

\textbf{Context:} an instructor records student attendance from the command
line.

\textbf{Instructions:}
\begin{enumerate}
  \item Create a script named \texttt{attendance.sh} that:
  \begin{itemize}
    \item Accepts student name and status arguments
      (\texttt{present}/\texttt{excused}/\texttt{absent})
    \item Saves each entry to \texttt{attendance-YYYY-MM-DD.txt} in the format
      \texttt{[HH:MM] NAME - STATUS}
    \item Supports the \texttt{-r} option to show a recap (count per status)
    \item Supports the \texttt{-h} option to show help
  \end{itemize}
  \item Record at least 5 entries and display the recap.
\end{enumerate}

\textbf{Required concepts:} variables, positional parameters,
\texttt{getopts}, \texttt{if}, \texttt{for}, functions, and file
redirection.

\subsection*{Assignment 2 System Health Check Script}

\textbf{Context:} an administrator performs a server health check before
maintenance.

\textbf{Instructions:}
\begin{enumerate}
  \item Create a script named \texttt{healthcheck.sh} using the professional
    template from the Best Practices section.
  \item The script should display: date/time, hostname, uptime, CPU usage,
    memory usage, and disk usage for every mounted filesystem.
  \item If any disk usage exceeds 80\%, display a warning.
  \item Save the result to \texttt{healthcheck-YYYY-MM-DD.log} and also show
    it in the terminal using \texttt{tee}.
  \item The option \texttt{-t <percent>} changes the disk warning threshold
    (default 80).
\end{enumerate}

\textbf{Required concepts:} \texttt{set -euo pipefail}, \texttt{trap},
\texttt{getopts}, functions with \texttt{local}, \texttt{for},
\texttt{if}, and \texttt{tee}.

%%---------------------------------------------------------------------
\section{Summary}
\label{sec:summary-bash-prog}

\begin{itemize}
  \item The shebang (\texttt{\#!/bin/bash}) determines the interpreter.
    Without a shebang, script behavior is undefined.
  \item Bash variables are untyped; arithmetic uses \texttt{\$(( ))}; command
    output is captured with \texttt{\$( )}.
  \item Positional parameters (\texttt{\$1}, \texttt{\$\#}, \texttt{\$@})
    let scripts receive arguments from the command line.
  \item The control structures \texttt{if}, \texttt{for}, \texttt{while},
    and \texttt{case} are conceptually equivalent to Java, with different
    syntax.
  \item Functions receive arguments via \texttt{\$1}/\texttt{\$2}; use
    \texttt{local} for local variables and \texttt{echo} to return values.
  \item \texttt{getopts} processes options like \texttt{-x} in standard Unix
    commands.
  \item \texttt{set -euo pipefail} and \texttt{trap} are the foundation of
    reliable, debuggable scripts.
\end{itemize}

%
